\section{简介}

本需求规格说明书用于说明“轻量级外卖服务平台”的需求来源、业务边界、功能行为和质量要求，并作为需求评审、系统设计、开发测试、交叉验收和版本变更的共同依据。文档面向课程教师、验收人员、项目经理、开发人员和测试人员。

\subsection{背景}

本项目源于《软件工程综合实践》课程的外卖业务实践任务。平台围绕顾客、商家、骑手和管理员四类业务参与者展开：顾客浏览商品并下单，商家维护商品并履行订单，骑手承接配送任务，管理员负责准入审核和异常治理。在完成基础交易闭环的同时，系统加入 AI 客服与智能点餐、营销资产、评价和数据看板等增强能力，但不让任何外部智能服务成为基础下单和配送的单点依赖。

\subsection{定义、缩略语}

\begin{longtable}{|L{0.22\linewidth}|L{0.68\linewidth}|}
\caption{术语与缩略语}\label{tab:terms}\\ \hline
\textbf{术语} & \textbf{定义} \\ \hline
有效订单 & 已支付且未取消的订单；营业额与销量默认仅统计已完成的有效订单。 \\ \hline
配送任务 & 订单进入配送阶段后生成的独立履约记录，记录骑手、任务状态和关键时间。 \\ \hline
活动任务 & 处于待接单、已接单、已取餐、配送中或异常待处理状态的配送任务。 \\ \hline
资产流水 & 钱包、积分和优惠券因业务动作产生的不可直接覆盖的变更记录。 \\ \hline
幂等 & 同一业务请求重复提交时，系统不产生重复扣款、扣库存、返还或状态跳转。 \\ \hline
AI 工具调用 & AI 客服在用户授权与后端权限校验后，对菜单搜索、订单查询等受控接口的调用。 \\ \hline
降级 & 外部 AI、地图或媒体服务不可用时，返回明确提示并切换到普通搜索、文本地址或手工操作。 \\ \hline
REST & Representational State Transfer，本项目后端接口采用的资源化交互风格。 \\ \hline
\end{longtable}

\subsection{约束}

\begin{longtable}{|L{0.20\linewidth}|L{0.70\linewidth}|}
\caption{项目约束}\label{tab:constraints}\\ \hline
\textbf{类别} & \textbf{约束内容} \\ \hline
课程过程 & 需求、设计、编码、测试和验收材料应保持可追踪；需求变更须同步更新文档和测试。 \\ \hline
技术环境 & Web 前端采用 Vue 3，后端采用 Spring Boot，持久化数据库采用 MySQL；前后端通过 REST 接口交互。 \\ \hline
业务边界 & 系统只模拟支付与钱包结算，不接入真实微信、支付宝或银行卡清算；本期不实现复杂自动派单和骑手薪资结算。 \\ \hline
数据权限 & 所有私有资源均由后端依据认证身份、角色和资源归属进行校验，不信任客户端传入的用户身份。 \\ \hline
外部服务 & AI、地图和对象存储通过后端适配接口接入；凭证仅保存在后端配置中，外部服务失败不得阻断核心交易。 \\ \hline
验收环境 & 使用项目提供的演示账号和初始化数据；在浏览器能够访问前后端服务时，应能完成 P0/P1 基线场景。 \\ \hline
\end{longtable}

\subsection{参考资料}

\begin{longtable}{|L{0.34\linewidth}|L{0.18\linewidth}|L{0.38\linewidth}|}
\caption{参考资料}\label{tab:references}\\ \hline
\textbf{资料名称} & \textbf{提供方/作者} & \textbf{用途} \\ \hline
《软件工程综合实践》课程资料与项目任务书 & 课程组 & 确定实践范围、过程要求和验收要求。 \\ \hline
《UML 大战需求分析》附录 1 & 张传波 & 作为本需求规格说明书的章节结构与用例描述参考。 \\ \hline
轻量级外卖服务平台公开仓库 & 项目组 & 保存对外源代码、需求、接口、测试和验收材料。 \\ \hline
\end{longtable}
